% Created 2026-03-12 Thu 17:25
% Intended LaTeX compiler: pdflatex
\documentclass[hidelinks, 12pt]{article}

\usepackage[utf8]{inputenc}
\usepackage[T1]{fontenc}
\usepackage{graphicx}
\usepackage{longtable}
\usepackage{wrapfig}
\usepackage{rotating}
\usepackage[normalem]{ulem}
\usepackage{amsmath}
\usepackage{amssymb}
\usepackage{capt-of}
\usepackage{hyperref}
\makeatletter \@ifpackageloaded{geometry}{\geometry{margin=2cm}}{\usepackage[margin=2cm]{geometry}} \makeatother
\usepackage{siunitx}
\usepackage{amsmath}
\usepackage{amssymb}
\usepackage{mathtools}
\usepackage{circuitikz}
\usepackage{float}
\usepackage{karnaugh-map}
\usepackage{libertine}
\usepackage[libertine,cmintegrals,cmbraces,vvarbb]{newtxmath}
\usepackage[scaled=0.95]{zi4}
\usepackage{fancyhdr}
\pagestyle{fancyplain}
\lhead{Bach \textsc{Wongwandanee}}
\rhead{ECE 403}
\chead{Assignment 8}
\author{Waridh (Bach) Wongwandanee}
\date{\today}
\title{ECE 403 Assignment 8}
\hypersetup{
 pdfauthor={Waridh (Bach) Wongwandanee},
 pdftitle={},
 pdfkeywords={},
 pdfsubject={},
 pdfcreator={},
 pdflang={English}}
\usepackage{biblatex}

\begin{document}

\section*{Question 1}
\label{sec:orgd8b70e1}
Find the fastest functional clock frequency for a 3-cycle pipeline that feeds back on itself.
Solve with each of 2-phase latches and pulse latches.
Verify that hold time is met. Use no time borrowing and the maximum time borrowing available.
\begin{itemize}
\item \(T_{skew} = \qty{50}{\pico\s}\), direction unknown
\item 7 groups of combinational logic, each with \(T_{cd\_group} = \qty{100}{\pico\s}\) and, used in the provided sequence, \(T_{pd} = \left\{ 500, 600, 300, 400, 500, 200, 110 \right\} \unit{\pico\s}\) Synchronize the \qty{110}{\pico\s} stage output to the clock.
\item latches and D-FFs : \(T_{setup} = \qty{40}{\pico\s}, T_{hold} = \qty{30}{\pico\s}, T_{pcq}=\qty{45}{\pico\s}, T_{ccq}=\qty{30}{\pico\s}, T_{pdq}=\qty{50}{\pico\s}, T_{cdq}=\qty{35}{\pico\s}\)
\end{itemize}
\subsection*{2 phase latch}
\label{sec:org7e0fad8}
In this section, for the 2-phase latch systems, we are going to start the problem with a \(t_{nonoverlap}=\qty{0}{\pico\s}\).
\subsubsection*{Without time borrowing}
\label{sec:org4fc234a}
For a two phase latch, we would actually split the 3-cycles into 6 stages.
They shall be grouped as follows.

\begin{verbatim}
{500, 600, 300, 400, 500, 200 + 110}
\end{verbatim}

Here, we are using the standard expression for setup time and hold time.
Additionally, we are assuming that the size of each half-stage shall be the same, which means that we will pick the larger of the two stages when trying to compute the minimum period.

\begin{equation}
\label{eq-2-phase-prop}
T_c \geq  t_{pdq1} + t_{pd1} + t_{pdq2} + t_{pd2} + t_{skew}
\end{equation}

\begin{equation}
\label{eq-2-phase-cont}
t_{cd1}, t_{cd2} \geq t_{hold} - t_{ccq} - t_{nonoverlap}
\end{equation}


\begin{table}[htbp]
\caption{\label{tab-q1-1}The computed values. The value of minimum \(T_{c}\) was computed from (\ref{eq-2-phase-prop}). The verification of the hold time is done through (\ref{eq-2-phase-cont}). In this case, \(t_{skew}=\qty{50}{\pico\s}\). Additionally, \(t_{cd}\) must be higher than \texttt{0} \unit{\pico\s}. The frequency was found to be \qty{740.74}{\mega\hertz}.}
\centering
\begin{tabular}{rrrrrrrl}
Stage & \(t_{pd1}\) & \(t_{pd2}\) & max \(t_{pd}\) & \(t_{cd1}\) & \(t_{cd2}\) & min \(T_{c}\) & meets hold time\\
\hline
1 & 500 & 600 & 600 & 100 & 100 & 1350 & yes\\
2 & 300 & 400 & 400 & 100 & 100 & 950 & yes\\
3 & 500 & 310 & 500 & 100 & 200 & 1150 & yes\\
\hline
Max period & 1350 &  &  &  &  &  & \\
Max frequency & 740740740.7407408 &  &  &  &  &  & \\
\end{tabular}
\end{table}
\subsubsection*{With time borrowing}
\label{sec:orgbf4d18b}
With maximum time borrowing, the original grouping of combinational logic would not get a large performance improvement.
We could rearrange the grouping as follows, and then try to improve the performance.

\begin{verbatim}
{600, 500, 400, 300, 500, 200 + 110}
\end{verbatim}

Now, the first group could borrow time from the second group.

\begin{equation}
\label{eq-borrow-time}
t_{borrow} = \frac{T_{c borrow}}{2} - \left( t_{setup} + t_{nonoverlap} + t_{skew} \right)
\end{equation}

So here is how we are going to approach finding the minimum \(T_{c}\). We are going to make a theoretical \(T_{c}\) so that we could compute the possible borrow value.
We are also computing the \(t_{borrow}\) that we actually need.
From there, if the amount required is more than the amount that we need, we are going to add the required amount to that period.
If we could borrow the amount, then we are at the optimal period length.

\begin{table}[htbp]
\caption{\label{tab-q1-2}The computed values. The value of minimum \(T_{c}\) was computed from (\ref{eq-2-phase-prop}) and (\ref{eq-borrow-time}). The verification of the hold time is done through (\ref{eq-2-phase-cont}). In this case, \(t_{skew}=\qty{50}{\pico\s}\). Additionally, \(t_{cd}\) must be higher than \texttt{0} \unit{\pico\s}. The frequency was found to be \qty{760.456}{\mega\hertz}.}
\centering
\begin{tabular}{rrrrrrrrrl}
Stage & \(t_{pd1}\) & \(t_{pd2}\) & \(T_{ctheo}\) & \(t_{borrowreq}\) & \(t_{borrowpos}\) & min \(T_{c}\) & \(t_{cd1}\) & \(t_{cd2}\) & valid \(t_{hold}\)\\
\hline
1 & 600 & 500 & 1250 & 100 & 35 & 1315 & 100 & 100 & yes\\
2 & 400 & 300 & 850 & 100 & 35 & 915 & 100 & 100 & yes\\
3 & 500 & 310 & 960 & 190 & 80 & 1070 & 100 & 200 & yes\\
\hline
Max period & 1315 &  &  &  &  &  &  &  & \\
Max freq & 7.60456e+08 &  &  &  &  &  &  &  & \\
\end{tabular}
\end{table}
\subsection*{Pulse latch}
\label{sec:orgefd6357}
For the pulsed latch case, we are going to setup the combinational stages similarly to the flip-flop implementation.

\begin{verbatim}
{500 + 600, 300 + 400, 500 + 200 + 110}
\end{verbatim}
\subsubsection*{Without time borrowing}
\label{sec:org83514eb}

\begin{equation}
\label{eq-pulse-prop}
T_c \geq \text{max} \left( t_{pdq} + t_{pd} , t_{pdq} + t_{pd} + t_{setup} - t_{pw}  \right) + t_{skew}
\end{equation}

\begin{equation}
\label{eq-pulse-cont}
t_{hold} + t_{pw} + t_{skew} \leq t_{ccq} + t_{cd}
\end{equation}

\begin{table}[htbp]
\caption{\label{tab-q2-1}The computed values. The value of minimum \(T_{c}\) was computed from (\ref{eq-pulse-prop}). The verification of the hold time is done through (\ref{eq-pulse-cont}). In this case, \(t_{skew}=\qty{50}{\pico\s}\). Here, we found the max \(t_{pw}\) that we could use in the design by rearranging (\ref{eq-pulse-cont}), which we found to be \qty{150}{\pico\s}. Additionally, the contamination delay of the combinational segment must be higher than \texttt{200} \unit{\pico\s}. The frequency was found to be \qty{833.33}{\mega\hertz}.}
\centering
\begin{tabular}{rrrrrl}
Stage & \(t_{pd}\) & total \(t_{cd}\) & max \(t_{pw}\) & min \(T_{c}\) & meets hold time\\
\hline
1 & 1100 & 200 & 150 & 1200 & yes\\
2 & 700 & 200 & 150 & 800 & yes\\
3 & 800 & 300 & 150 & 900 & yes\\
\hline
Max period & 1200 &  &  &  & \\
Max freq (\unit{\hertz}) & 8.33333e+08 &  &  &  & \\
\end{tabular}
\end{table}
\subsubsection*{With time borrowing}
\label{sec:orgc162f49}
We don't have to change any ordering of the combinational logic here, we are just trying to borrow time from the pulse.
We are going to use the previously derived value for the pulse with to see how much time we could borrow, and then compute the new period.

\begin{equation}
\label{eq-pulse-borrow}
t_{borrow} \leq t_{pw} - \left( t_{setup} + t_{skew} \right)
\end{equation}

\begin{table}[htbp]
\caption{\label{tab-q2-2}The computed values. The value of minimum \(T_{c}\) was computed from (\ref{eq-pulse-prop}), as well as some borrowing from (\ref{eq-pulse-borrow}). The verification of the hold time is done through (\ref{eq-pulse-cont}). In this case, \(t_{skew}=\qty{50}{\pico\s}\). Here, we found the max \(t_{pw}\) that we could use in the design by rearranging (\ref{eq-pulse-cont}), which we found to be \qty{150}{\pico\s}. Additionally, the contamination delay of the combinational segment must be higher than \texttt{200} \unit{\pico\s}. The frequency was found to be \qty{877.19}{\mega\hertz}.}
\centering
\begin{tabular}{rrrrrrl}
Stage & \(t_{pd}\) & total \(t_{cd}\) & max \(t_{pw}\) & \(t_{borrow}\) & min \(T_{c}\) & meets hold time\\
\hline
1 & 1100 & 200 & 150 & 60 & 1140 & yes\\
2 & 700 & 200 & 150 & 60 & 740 & yes\\
3 & 800 & 300 & 150 & 60 & 840 & yes\\
\hline
Max period & 1140 &  &  &  &  & \\
Max freq (\unit{\hertz}) & 8.77193e+08 &  &  &  &  & \\
\end{tabular}
\end{table}
\end{document}
